\vspace{-2ex}
\section{Introduction} 
\label{sec:introduction}

Blockchain technology has revolutionized the creation of global, secure, and
programmable ledgers, fundamentally altering how digital transactions are
conducted. Central to this innovation are smart contracts, which operate on
blockchains, allowing developers to define and enforce complex
transactional rules directly on the ledger. This capability has positioned
smart contracts as the backbone of various decentralized
financial (DeFi) services. As of March 13th, 2025, the total value locked in
3,973 DeFi protocols has surged to approximately \$87.82
billion~\cite{defillama}, highlighting the substantial economic impact and
growth of this technology.

However, by the same date, vulnerabilities in DeFi smart contracts have resulted in 
financial losses exceeding \$11.21 billion USD~\cite{defillamahacks}. 
In response, researchers have developed program analysis and verification techniques 
to secure smart contracts~\cite{luu2016making, torres2018osiris, 
liu2018reguard, feist2019slither, zhang2019mpro, bose2022sailfish, 
ma2021pluto, zheng2022park, wang2024efficiently, albert2020taming, 
zhang2023your}, and developers often commission security audits prior to deployment.


% However, the increasing reliance on smart contracts has brought security
% concerns to the forefront. By the same date, vulnerabilities in DeFi smart
% contracts have resulted in financial losses exceeding \$11.21 billion
% USD~\cite{defillamahacks}. In response to these growing threats, researchers have
% developed an array of program analysis and verification techniques and tools aimed at
% securing smart contracts~\cite{luu2016making, torres2018osiris, liu2018reguard,
% feist2019slither, zhang2019mpro, bose2022sailfish, ma2021pluto, zheng2022park,
% wang2024efficiently, albert2020taming, zhang2023your}. 
% Additionally, to mitigate risks and ensure the
% integrity of smart contracts, developers often commission thorough security
% audits to identify and address potential errors and vulnerabilities prior to the
% contract deployment.



Despite these advancements in security measures, the evolving landscape of smart 
contracts has continually outpaced traditional defenses.
Modern smart
contracts are now designed with the flexibility to support the layering of
additional contracts, a feature particularly vital in DeFi ecosystem. In DeFi, smart contracts facilitate a diverse array of
financial products and services, such as lending and
yield farming. These interlinked contracts are often referred to as ``DeFi
legos,'' emphasizing their modularity. The ability to combine various DeFi smart
contracts, a concept known as ``DeFi composability,'' is widely regarded as one
of the key advantages of DeFi~\cite{popescu2020decentralized, amler2021defi,
schar2021decentralized}. 
However, this complexity and interdependence introduce
significant challenges to securing smart contracts with conventional methods.
The security of a DeFi protocol depends not only on the correct design and 
implementation of its own contracts but also on the integrity of external 
contracts it interacts with. Additionally, experienced users or attackers can 
deploy their own smart contracts to invoke functions across 
multiple DeFi protocols in arbitrary sequences. 
Considering all possible interactions a DeFi 
protocol may encounter prior to deployment is infeasible for traditional 
security approaches.


%This interconnectedness exacerbates the difficulty of addressing well-known
%vulnerabilities. \chen{ the third sentence doesn't go well with the previous 
%two sentences, since simple re-entrancy like DAO 2016 attack wasn't not caused by 
%the interconnectedness. So, I think it is better to remove this sentence since 
%reviewers might assume that we think DAO 2016 attack is related to 
%interconnectedness and assume we don't know what we are talking about.
%} For instance, the re-entrancy vulnerability, infamously
%exploited in the 2016 DAO hack, allowed a hacker to repeatedly invoke a
%vulnerable contract function~\cite{liu2018reguard, coindesk_dao_hack}. Despite
%heightened awareness of this kind of vulnerabilities, attackers have since
%developed a variety of new sophisticated
%techniques~\cite{caversaccio_reentrancy_attacks} to exploit smart contracts,
%including read-only re-entrancy, cross-function re-entrancy, cross-contract
%re-entrancy, and cross-chain re-entrancy~\cite{callens2024temporarily,
%valixconsulting_cross_function_reentrancy_2023, wang2024efficiently,
%valixconsulting_cross_contract_reentrancy_2023, cross_chain_reentrancy}. These
%evolving attack methods highlight the inherent limitations of current security
%measures.

\if 0
Yet, with the increasing significance of smart contracts in the blockchain
ecosystem, their security has emerged as a pivotal issue. As of August 14th,
2024, vulnerabilities in DeFi smart contracts have led to financial damages
exceeding \$6.05 billion USD~\cite{defillamahacks}. Consequently, reinforcing
the reliability and security of smart contracts is crucial to maintaining the
integrity of decentralized applications.

As demands for more functionalities increase, smart contracts have become
significantly more complex. Additionally, they are designed with the
flexibility to support the building of additional contracts on top of one
another. This is particularly crucial in the DeFi ecosystem, where smart
contracts enable a wide range of financial products and services such as
lending, borrowing, trading, and yield farming. These smart contracts are
sometimes referred to as ``DeFi legos''. The ability to compose different DeFi
smart contracts, known as ``DeFi composability,'' is often seen as one of the
main benefits of DeFi~\cite{popescu2020decentralized, amler2021defi,
schar2021decentralized}.

However, this flexibility also introduces security vulnerabilities. While the
composability of DeFi smart contracts enables ethical developers to craft
innovative applications atop existing contracts, it also opens the door for
malicious actors to exploit unforeseen uses of these contracts to orchestrate
exploits in ways that developers did not predict. 

Current approaches to securing smart contracts are only partially effective
against unforeseen exploits. Developers typically commission security audits to
identify vulnerabilities, but these audits are costly and time-consuming.
Moreover, they may not detect all potential threats. When alerted to new
threats, developers can also promptly inspect their smart contracts for
weaknesses and update them using proxies~\cite{bodell2023proxy}. However, this
reactive strategy requires constant vigilance regarding emerging security
breaches, and developers often find it challenging to outpace hackers in
identifying and patching vulnerabilities. A notable example is the re-entrancy
flaw, infamously exploited during the DAO hack of 2016, where a hacker managed
to re-entrant a single function of the vulnerable
contract~\cite{liu2018reguard, coindesk_dao_hack}. Despite increased awareness
of this flaw, attackers have developed multiple new methods to exploit many
protocols~\cite{caversaccio_reentrancy_attacks}, including read-only
re-entrancy, cross-function re-entrancy, cross-contract re-entrancy, and
cross-chain re-entrancy~\cite{callens2024temporarily,
valixconsulting_cross_function_reentrancy_2023, wang2024efficiently,
valixconsulting_cross_contract_reentrancy_2023, cross_chain_reentrancy}. These
innovations underscore the limitations of current security measures. This
ongoing challenge highlights the inadequacy of existing security protocols in
adapting to the evolving threat landscape with unforeseen exploits.
\fi

A critical observation in hack transaction analysis is that they often
\emph{exploit unintended control flows across multiple functions}, deviating
from the original design intentions of the developers. For instance,
re-entrancy attacks leverage an unforeseen recursive control flow through
default handlers in custom contracts, enabling repeated execution of a critical
function within one transaction. Similarly, flash loan attacks manipulate
multiple functions in a precisely timed sequence, utilizing large asset
transfers to coerce the victim contract into executing unfavorable trades.
We conducted an empirical study analyzing transaction histories of 
\revise{$37$} compromised Ethereum protocols. 
Our findings reveal that control flows are relatively constrained, 
with attack transactions introducing novel flows never observed 
previously in all but one hacked protocol.


% To further explore this phenomenon, we conducted an empirical study analyzing
% transaction histories of \revise{$37$} compromised protocols on
% Ethereum. Our findings reveal that the number of different control flows is
% relatively constrained. In all but one cases, attack transactions
% introduced novel control flows that had never been observed in any previous
% transaction.

%Based on our findings, we developed a control flow integrity technique, {\name}, that combines
%dynamic and static analysis to track and whitelist control flows in smart
%contracts. {\name} consists of simple control flow whitelisting policies and a set of control flow
%reduction rules. It captures control flows and runtime behaviors, enforcing
%these policies and rules on the fly. 

\noindent \textbf{{\name}:} 
Building on the above observations, we developed {\name}, a novel framework to
enforce \emph{control flow integrity} to secure smart contracts. Given a DeFi
protocol, {\name} instruments its existing smart contracts with additional code
to track control flow data. {\name} also deploys a new guard contract that
collects control flow data from these instrumented contracts, and enforces four
whitelisting policies at runtime to detect and neutralize any attacks that
attempt to exploit unexpected control flows. Unlike many previous invariant
enforcement tools~\cite{chen2024demystifying, liu2022invcon}, {\name} does not
rely on inferring its security rules from prior benign transaction traces and
therefore can apply to smart contracts immediately at their initial
deployments, leaving no gap of unprotected periods.

% Note that control flow integrity is a well-established security technique in
% traditional software to prevent memory attacks from hijacking a program's
% control flow~\cite{burow2017control, abadi2009control, goktas2014out}. Though
% our proposed techniques in this paper are dramatically different from prior
% control flow integrity techniques for traditional software, we adapt this
% concept here to emphasize the critical role of control flow in our approach.


% One of the primary challenges for {\name} is to accurately detect potentially
% malicious control flows while minimizing false positives. Although the number
% of unique control flows is inherently limited, a naive whitelisting approach
% could still produce numerous false positives or demand substantial human
% intervention. To overcome this, {\name} leverages read-after-write dependency
% analysis to simplify control flow traces. It does so by recording the ledger
% state accessed by each function call and monitoring the dependencies between
% these calls. When function calls within a trace are independent—that is, they
% do not rely on one another--{\name} treats them as separate control flows and
% applies whitelisting policies individually. This strategy effectively
% eliminates redundant traces and significantly lowers the false positive rate.

% A key challenge arises from {\name}'s whitelisting-only design. 
% To prevent false positives, {\name} employs heuristics to simplify control flows: 
% excluding read-only calls that don't alter state, 
% tracking read-after-write dependencies, and treating independent calls as 
% separate flows. These heuristics effectively reduce redundant traces 
% and lower false positive rates.



% A key challenge for {\name} arises from its fundamental design to only
% whitelist control flows rather than using prior knowledge of hacks to blacklist
% potentially malicious ones. 

A key challenge arises from {\name}'s whitelisting-only design. 
This approach, while enhancing security by adhering
strictly to known safe paths, could inherently lead to an increase in false
positives if not meticulously managed. Although the number of unique control
flows is inherently limited, a naive whitelisting approach could still produce
numerous false positives or demand substantial human intervention. To address
this issue,  {\name} employs heuristics to simplify control flows: 
excluding read-only calls that don't alter state, 
tracking read-after-write dependencies, and treating independent calls as 
separate flows. These heuristics effectively simplify control flows,
and lower false positive rates.

% Specifically, {\name} excludes all read-only function
% calls from the control flows, as these calls do not alter the blockchain
% ledger's state. Additionally, {\name} records the ledger state accessed by each
% function call and tracks the read-after-write dependencies of these calls. If
% function calls in a control flow trace do not have any dependencies with each
% other, {\name} will treat them as separate control flows and assess them with
% whitelisting policies individually. These heuristics effectively reduce
% redundant traces and significantly lower the false positive rate for {\name}. 

\noindent \textbf{Experimental Results:}
We evaluated {\name} on the deployed smart contracts and their transactions 
of the \revise{$37$} hacked DeFi protocols
included in our empirical study. Our results indicate 
that, when 
configured only once at deployment, {\name} can effectively prevent 
\revise{$35$} out of \revise{$37$}
attacks analyzed in our study, maintaining a low average false positive rate of
just \revise{0.26\%}. Unlike traditional methods, {\name} does not depend on historical
transactions. Despite this, {\name} still surpasses the state-of-the-art which 
instruments the smart contracts with invariants learned from historical transactions.
Moreover, after implementing \revise{four}
optimization techniques, {\name} achieves a minimal gas consumption overhead of
\revise{3.53\%} on average. These results demonstrate the usefulness of our empirical
findings and {\name}. 

\noindent\textbf{Contributions:} This paper makes the following contributions. 
% \noindent \textbf{Contributions:} To summarize, we make the following
% contributions in this paper: 
\begin{itemize}[leftmargin=*]
    \item \textbf{Empirical Study:} To the best of our knowledge, we conducted
        the first comprehensive empirical study of control flows in historical
        transactions of compromised DeFi protocols. Our analysis uncovers
        critical insights into the control flow patterns prevalent in DeFi protocols and
        explores various use cases of DeFi composability.

    \item \textbf{{\name} Technique:} This paper proposes the first control flow integrity
        technique for smart contracts with whitelisting policies and
        simplification heuristics. This paper also details methods for implementing 
        these policies and heuristics through static and dynamic analysis, and describes
        how they are instrumented in contracts and enforced on the fly. 

    \item \textbf{Evaluation and Tools:} This paper evaluates the effectiveness
        of {\name} in preventing attacks. To support ongoing research and
        facilitate community engagement, we provide open access to the
        study results, experimental results, and our tool,
        available at our website~\cite{website}.
\end{itemize}
\vspace{-2mm}


